\chapter{Introduction}

Programs for machine learning, simulation, optimization, and dynamic programming
are often written as tensor equations. The programmer names axes, sums over
intermediate dimensions, carries state through a recurrence, and eventually asks
how one quantity changes with respect to another. Those activities are
mathematically related, but contemporary programming systems often expose them
through separate interfaces.

For example, a computation may use an \texttt{einsum} string to express a
contraction, a host-language loop or scan combinator to express recurrence, and
a separate automatic-differentiation transform to obtain a gradient. The program
executes, but the compiler-facing facts are scattered: output axes are hidden in
one notation, recurrence dependencies in another, and derivative targets in a
third. The central question of this thesis is whether those facts can remain in
one source language without sacrificing executable behavior.

\section{Thesis Statement}

The thesis is that indexed tensor definitions, declarative recurrence, and local
derivative requests form a productive abstraction boundary for tensor-heavy
programs. When these constructs share one binding environment, the source
program keeps mathematical structure visible long enough for compiler passes to
check shape consistency, derive recurrence order, choose storage windows, and
lower derivative requests as shaped values.

\section{Implementation Scope}

This document is intentionally broader than the paper. The paper isolates the
language-design claim. The thesis treats \einlang{} as an implementation
artifact and documents the substantial subsystems in the repository:

\begin{table}[htbp]
\centering
\caption{Implementation subsystems covered by this thesis.}
\label{tab:implementation-map}
\renewcommand{\arraystretch}{1.08}
\begin{tabular}{@{}L{0.32\textwidth}L{0.56\textwidth}@{}}
\toprule
Subsystem & Non-trivial responsibility \\
\midrule
Frontend & Lark grammar, parse caching, source locations, AST construction, string interpolation, functions, literals, patterns, and Einstein syntax. \\
Identity and modules & DefId allocation, scoped resolution, Rust-like module paths, stdlib discovery, imports, aliases, re-exports, and Python module hooks. \\
IR and passes & Slot-based IR nodes, visitors, S-expression serialization, pass ordering, \texttt{TyCtxt}, range, shape, type, validation, pruning, lowering, and tree shaking. \\
Tensor lowering & Einstein declaration grouping, rest-index preprocessing, implicit range discovery, lowerings for reductions, selections, comprehensions, and backend execution facts. \\
Recurrence & Dependency extraction, recurrence-dimension ordering, same-timestep dependencies, storage-window metadata, and circular-buffer runtime execution. \\
Autodiff & Source \texttt{@} forms, differential IR, high-level differentiation, request lowering, leak checks, runtime JVP/VJP templates, and lazy Jacobian values. \\
Runtime and backends & DefId-keyed execution environment, NumPy expression and Einstein interpreters, vectorization, profiling/debug hooks, IREE lowering subset, and fallback behavior. \\
Library and evidence & Standard library modules, examples, reference implementations, unit/integration/stdlib/example tests, and docs contracts. \\
\bottomrule
\end{tabular}
\end{table}

\section{Non-Focuses}

It is equally important to state what this implementation is not. \einlang{} is
not currently a bytecode VM, a garbage-collected native runtime, or a compiler
whose central technical step is conversion from mutable code into SSA. The
repository contains an interpreter-style NumPy backend and an experimental IREE
lowering path; Python and NumPy own ordinary object allocation and memory
management.

\begin{table}[htbp]
\centering
\caption{Implementation areas that are not thesis claims.}
\label{tab:non-focus}
\renewcommand{\arraystretch}{1.08}
\begin{tabular}{@{}L{0.34\textwidth}L{0.54\textwidth}@{}}
\toprule
Not claimed & Reason \\
\midrule
Mutable-to-SSA conversion & Source bindings are immutable, but the compiler IR is a mutable, source-located tree annotated by passes; there is no CFG-level SSA conversion or phi-node pass. \\
Bytecode generation and register allocation & The main backend executes IR through Python/NumPy visitors. The IREE backend emits MLIR for a limited function subset, not an Einlang bytecode format. \\
Custom virtual machine & There is no register VM, threaded dispatch loop, or bytecode interpreter in the current runtime. \\
Custom garbage collector & Memory is managed by Python and NumPy. Recurrence circular buffers are storage optimizations, not a GC. \\
C FFI and pinning & The current implementation has Python module hooks and NumPy file I/O, but no C FFI with pinning or ownership transfer. \\
Parser generator as whole frontend & Lark supplies parsing machinery, but the grammar, transformers, source locations, AST normalization, diagnostics, and module integration are implementation work. \\
\bottomrule
\end{tabular}
\end{table}

\section{Contributions}

This work makes five contributions.
\begin{itemize}
\item It identifies indexed recurrence with local derivative requests as a useful
source boundary for programs that mix tensor algebra, dynamic programming,
sequence models, and sensitivity analysis.
\item It presents \einlang{}, a source language in which multidimensional
indexed definitions, recurrent clauses, and derivative requests share names and
scope.
\item It describes the compiler contracts that make the constructs
compositional: clause-scoped index domains, static self-read offsets, shaped
derivative values, and DefId-stable binding identity.
\item It gives concrete compiler and runtime analyses for range/shape
inference, Einstein lowering, recurrence order, recurrence storage, autodiff,
and backend execution.
\item It reports on a NumPy-backed implementation, an experimental IREE backend,
and a test/example suite that exercises the major language and implementation
paths.
\end{itemize}

\begin{figure}[p]
\centering
\resizebox{\textwidth}{!}{\begin{tikzpicture}[
  font=\sffamily\footnotesize,
  >=Latex,
  title/.style={font=\sffamily\bfseries\large,text=figInk,align=center},
  box/.style={draw=figLine,line width=0.6pt,rounded corners=2pt,fill=white,align=left,inner sep=7pt,text=figInk},
  fact/.style={box,minimum width=3.3cm,minimum height=1.25cm},
  consumer/.style={box,minimum width=3.3cm,minimum height=1.25cm,fill=figSlateBg},
  spine/.style={box,minimum width=11.4cm,minimum height=1.0cm,align=center},
  backend/.style={box,minimum width=5.0cm,minimum height=1.1cm,fill=white},
  arrow/.style={->,line width=0.8pt,draw=figLine},
  bus/.style={line width=0.8pt,draw=figLine}
]

\node[title] (title) at (0,0) {Source Facts and Their Implementation Consumers};

\node[spine,fill=figBlueBg,draw=figBlue] (source) at (0,-1.25) {
  \textbf{One Einlang source program}\\
  indexed bindings, recurrence clauses, derivative requests, modules, and local names
};

\node[fact,fill=figBlueBg,draw=figBlue] (axes) at (-4.1,-3.05) {
  \textbf{Axis facts}\\
  output indices\\
  reduction indices\\
  rest-pattern spans
};
\node[fact,fill=figGreenBg,draw=figGreen] (deps) at (0,-3.05) {
  \textbf{Dependency facts}\\
  self-read offsets\\
  same-timestep reads\\
  later observations
};
\node[fact,fill=figRoseBg,draw=figRose] (diffs) at (4.1,-3.05) {
  \textbf{Derivative facts}\\
  target bindings\\
  wrt bindings\\
  tangent axes
};

\node[consumer] (shape) at (-4.1,-5.35) {
  \textbf{Range, shape, type}\\
  infer domains\\
  check ranks\\
  type expressions
};
\node[consumer] (storage) at (0,-5.35) {
  \textbf{Recurrence planning}\\
  choose order\\
  prove windows\\
  annotate storage
};
\node[consumer] (ad) at (4.1,-5.35) {
  \textbf{Autodiff lowering}\\
  synthesize IR\\
  project Jacobians\\
  enforce leak checks
};

\node[spine,fill=figGoldBg,draw=figGold] (lowered) at (0,-7.65) {
  \textbf{Lowered, DefId-keyed implementation IR}\\
  explicit loops, guards, bindings, recurrence metadata, and backend execution facts
};

\node[backend] (numpy) at (-2.85,-9.75) {
  \textbf{NumPy backend}\\
  scalar, vectorized, hybrid\\
  reductions and circular buffers
};
\node[backend] (iree) at (2.85,-9.75) {
  \textbf{IREE subset}\\
  lazy MLIR emission\\
  signature cache and fallback
};

\draw[arrow] ([xshift=-4.1cm]source.south) -- (axes.north);
\draw[arrow] (source.south) -- (deps.north);
\draw[arrow] ([xshift=4.1cm]source.south) -- (diffs.north);

\draw[arrow] (axes.south) -- (shape.north);
\draw[arrow] (deps.south) -- (storage.north);
\draw[arrow] (diffs.south) -- (ad.north);

\coordinate (lowerbus) at (0,-6.55);
\draw[bus] (-4.1,-6.55) -- (4.1,-6.55);
\draw[bus] (shape.south) -- (-4.1,-6.55);
\draw[bus] (storage.south) -- (0,-6.55);
\draw[bus] (ad.south) -- (4.1,-6.55);
\draw[arrow] (lowerbus) -- (lowered.north);

\coordinate (backendbus) at (0,-8.75);
\draw[arrow] (lowered.south) -- (backendbus);
\draw[bus] (-2.85,-8.75) -- (2.85,-8.75);
\draw[arrow] (-2.85,-8.75) -- (numpy.north);
\draw[arrow] (2.85,-8.75) -- (iree.north);

\end{tikzpicture}
}
\caption{Thesis-specific implementation view: source facts are preserved until
the compiler pass that consumes them, then carried into lowered backend-facing
IR.}
\label{fig:fragmentation}
\end{figure}
